% Part: methods
% Chapter: proofs
% Section: using-definitions

\documentclass[../../../include/open-logic-section]{subfiles}

\begin{document}

\olfileid{mod}{prf}{def}

\olsection{运用定义}

我们曾提到，必须熟悉证明中可能用到的所有定义，并能正确地加以运用。这一点极其重要，值得稍作深入探讨。定义的作用是简写性质与关系，以便我们可以更简洁地谈论它们。引入的简写称为\emph{被定义项}，它所简写的部分称为\emph{定义项}。在证明中，我们经常需要回溯被定义项的引入方式，因为必须利用定义项的逻辑结构（被定义项正是其简写形式）来完成证明。通过展开定义，就能确保自己触及了逻辑推导的关键所在。

我们从一个例子开始。假设你想证明以下命题：

\begin{prop}
对于任意集合 $A$ 和 $B$，$A \cup B = B \cup A$。
\end{prop}

为了着手证明，我们首先需要知道两个集合“相等”意味着什么；也就是说，我们需要知道该等式中的“$=$”对集合而言是什么意思。集合被定义为相等，当且仅当它们具有相同的元素。因此，我们需要展开的定义是：

\begin{defn}
集合 $A$ 和 $B$ 是\emph{相等的}，记作 $A = B$，当且仅当 $A$ 的每个元素都是 $B$ 的元素，反之亦然。
\end{defn}

这个定义使用 $A$ 和 $B$ 作为任意集合的占位符。它所定义的——即\emph{被定义项}——是表达式“$A = B$”，并给出了 $A = B$ 为真的条件。这个条件——“$A$ 的每个元素都是 $B$ 的元素，反之亦然”——就是\emph{定义项}。\footnote{在这个特定情况下——而且非常令人困惑！——当 $A = B$ 时，集合 $A$ 和 $B$ 其实是同一个集合，尽管我们在等号左右用了不同的字母来表示它。但是，指称该集合的方式可能不同，这使得该定义并非平庸。}该定义指明，$A = B$ 为真当且仅当（我们将其缩写为“iff”）该条件成立。

当你应用这个定义时，必须将定义中的 $A$ 和 $B$ 匹配到你正在处理的情形。在我们的例子中，这意味着要使 $A \cup B = B \cup A$ 为真，每个 $z \in A \cup B$ 也必须在 $B \cup A$ 中，反之亦然。命题中的表达式 $A \cup B$ 扮演了定义中 $A$ 的角色，而 $B \cup A$ 则扮演了 $B$ 的角色。由于 $A$ 和 $B$ 同时在定义和我们要证明的命题陈述中被使用，但用法不同，你必须小心确保不会将两者混淆。例如，如果认为可以通过证明 $A$ 的每个元素都是 $B$ 的元素（反之亦然）来证明该命题，那就错了——那将证明的是 $A = B$，而不是 $A \cup B = B \cup A$。（此外，因为 $A$ 和 $B$ 可以是任意两个集合，这样做也不会走多远：如果对 $A$ 和 $B$ 没有任何假设，它们很可能就是不同的集合。）

在证明中，我们处理的是并集这样的集合论概念，因此还必须知道符号 $\cup$ 的含义，才能理解证明应如何推进。有时，展开一个定义会带来更多需要展开的定义。例如，$A \cup B$ 被定义为 $\Setabs{z}{z \in A
  \text{ or } z \in B}$。所以，如果想证明 $x \in A \cup B$，展开 $\cup$ 的定义可知，必须证明 $x \in \Setabs{z}{z \in A \text{ or } z \in B}$。现在你还需要记住，$x \in \Setabs{z}{\dots z\dots}$ 当且仅当 $\dots x\dots$。因此，进一步展开 $\Setabs{z}{\dots z \dots}$ 这种记法的定义，你需要证明的是：$x \in A$ 或 $x \in B$。所以，“$A \cup B$ 的每个元素也是 $B \cup A$ 的元素”实际上意味着：“对每个 $x$，如果 $x \in A$ 或 $x \in B$，那么 $x \in B$ 或 $x \in A$。”如果我们把命题中的定义完全展开，就会发现需要证明的是：

\begin{prop}
对于任意集合 $A$ 和 $B$：(a) 对每个 $x$，如果 $x \in A$ 或 $x \in B$，那么 $x \in B$ 或 $x \in A$；并且 (b) 对每个 $x$，如果 $x \in B$ 或 $x \in A$，那么 $x \in A$ 或 $x \in B$。
\end{prop}

重要的是，展开定义是构建证明的一个必要步骤。正确地做到这一点有时并不容易：你必须仔细区分并匹配定义中的变元与你所要证明的断言中的项。为了成功，你必须理解问题在问什么，以及问题中使用的所有术语分别是什么意思——你常常需要展开不止一个定义。在下面这样的简单证明中，解答几乎能直接从定义本身得出。当然，情况不会总是这么简单。

\begin{prob}
假设要求你证明 $A \cap B \neq \emptyset$。请展开其中出现的所有定义，即，用不提及“$\cap$”、“$=$”或“$\emptyset$”的方式重述这个断言。
\end{prob}

\end{document}